\textbf{1. Kiến thức cần nhớ}
\begin{itemize}
\item \textbf{HÌNH TRỤ}: Tạo thành khi quay hình chữ nhật một vòng quanh một cạnh cố định:
\begin{itemize}
\item Chiều cao $h$, bán kính đáy $R$, đường sinh $l$ ($l = h$).
\item \textbf{Diện tích xung quanh}: $\boxed{S_{xq} = 2\pi R h}$.
\item \textbf{Diện tích toàn phần}: $\boxed{S_{tp} = 2\pi R h + 2\pi R^2 = 2\pi R(h + R)}$.
\item \textbf{Thể tích}: $\boxed{V = S_{\text{đáy}} \cdot h = \pi R^2 h}$.
\end{itemize}
\item \textbf{HÌNH NÓN}: Tạo thành khi quay tam giác vuông một vòng quanh một cạnh góc vuông cố định:
\begin{itemize}
\item Chiều cao $h$, bán kính đáy $R$, đường sinh $l$.
\item Mối liên hệ Pythagore: $\boxed{l^2 = h^2 + R^2 \implies l = \sqrt{h^2 + R^2}}$.
\item \textbf{Diện tích xung quanh}: $\boxed{S_{xq} = \pi R l}$.
\item \textbf{Diện tích toàn phần}: $\boxed{S_{tp} = \pi R l + \pi R^2 = \pi R(l + R)}$.
\item \textbf{Thể tích}: $\boxed{V = \dfrac{1}{3} \pi R^2 h}$.
\end{itemize}
\end{itemize}

\textbf{2. Phương pháp giải các dạng}

%% DANG 9H0?31-1
\textbf{Dạng 1. Diện tích xung quanh, thể tích hình trụ}
\begin{enumerate}
\item Xác định bán kính đáy $R$ (nếu đề cho đường kính $d$ thì $R = \dfrac{d}{2}$) và chiều cao $h$.
\item Áp dụng công thức $S_{xq} = 2\pi R h$ và $V = \pi R^2 h$.
\item Chú ý đổi đơn vị thể tích: $1\text{ lít} = 1\text{ dm}^3 = 1000\text{ cm}^3 = 0{,}001\text{ m}^3$.
\end{enumerate}

%% DANG 9H0?31-2
\textbf{Dạng 2. Diện tích xung quanh, thể tích hình nón}
\begin{enumerate}
\item Xác định hai trong ba yếu tố ($R, h, l$), dùng $l^2 = h^2 + R^2$ để tìm yếu tố còn lại.
\item Áp dụng $S_{xq} = \pi R l$ và $V = \dfrac{1}{3}\pi R^2 h$.
\end{enumerate}

%% DANG 9H0?31-3
\textbf{Dạng 3. Bài toán thực tế về hình trụ, hình nón (Bài 5 đề Tuyển sinh 10)}
\begin{enumerate}
\item Nhận diện đồ vật: lon nước ngọt, thùng phi, ống thoát nước (hình trụ); nón lá, kem ốc quế, phễu, chóp nón giao thông (hình nón).
\item Phân biệt diện tích sơn/bọc: không đáy (chỉ tính $S_{xq}$), có $1$ đáy ($S_{xq} + \pi R^2$), có cả $2$ đáy ($S_{tp}$).
\item Dung tích chứa chất lỏng: tính thể tích $V$ rồi nhân với phần trăm chứa (ví dụ: nước chiếm $80\%$ thể tích).
\end{enumerate}
